\documentclass[11pt]{article}
\usepackage[a4paper,margin=2.2cm]{geometry}
\usepackage[T1]{fontenc}
\usepackage{mathpazo}
\usepackage{titlesec}
\usepackage{enumitem}
\usepackage{parskip}
\usepackage{xcolor}

\definecolor{konstanzblue}{HTML}{00539F}

\titleformat{\section}{\sffamily\bfseries\Large\color{konstanzblue}}{}{0em}{}
\titlespacing{\section}{0pt}{1.2em}{0.6em}

\newcommand{\entry}[3]{%
  \noindent\textbf{#1} \hfill \textit{#2}\\
  #3
  \vspace{0.9em}
  \par
}

\begin{document}
\begin{center}
  {\sffamily\bfseries\huge Trust and Credibility Section --- Source Summaries}\\[4pt]
  {\sffamily\large Chapter 2, \S 2.2 --- one to three sentences per source}
\end{center}
\vspace{1em}

\section*{Trust theory foundations}

\entry{Mayer, Davis \& Schoorman}{1995}{
Proposes an integrative model of organizational trust with three independently evaluable components of trustworthiness: ability, benevolence, and integrity. Defines trust itself as the willingness to be vulnerable to another party's actions, based on expected appropriate behaviour rather than the ability to monitor or control that party.
}

\entry{McKnight \& Chervany}{2001}{
Extends trust theory into a typology that separates trusting beliefs from trusting intentions. This distinction allows trust to be measured across both attitudinal (belief-based) and behavioural (intention-based) dimensions.
}

\entry{Alipour, Hartmann \& Alimardani}{2025}{
A PRISMA-compliant systematic review of 53 empirical studies at the intersection of uncanny valley effects and trust. Proposes a taxonomy of how trust has been operationalised in the literature and finds that most prior studies relied on static images or indirect encounters with self-report measures alone, recommending dynamic stimuli and behavioural indicators for future work.
}

\section*{Anthropomorphism, social presence, and credibility}

\entry{Nowak \& Biocca}{2003}{
Shows that anthropomorphism and perceived agency independently affect social presence and copresence. Very high levels of visual anthropomorphism do not consistently produce higher social presence, suggesting diminishing perceptual returns beyond a certain threshold.
}

\entry{Nowak \& Rauh}{2005}{
In a study with 255 participants rating a series of avatars, both perceived anthropomorphism and avatar gender expression predicted credibility judgments. More anthropomorphic avatars were rated as more credible, establishing that avatar design choices function as meaningful social signals.
}

\section*{Avatar stylisation and trust (direct comparisons)}

\entry{Canales et al.}{2024}{
Compared virtual doctors rendered at three stylisation levels (realistic, semi-realistic, caricatured) on multi-component trust before and after a short interaction, plus a behavioural forced-choice measure. Self-reported trust was statistically flat across style levels, but the semi-realistic avatar attracted the highest rate of behavioural selection.
}

\entry{Gao et al.}{2025}{
Compared realistic and stylised virtual agents across trust, social presence, and perceived empathy. Stylised agents achieved trust levels statistically comparable to realistic agents, voice characteristics had a negligible effect, and perceived agent age was a significant moderator of trust.
}

\entry{Alimardani et al.}{2024}{
A 2$\times$2 factorial study isolating visual appearance from voice quality. Mismatches between the two modalities elevated eeriness but produced no significant change in participants' rate of adopting the agent's recommendations, a behavioural trust measure.
}

\entry{Dubosc et al.}{2025}{
A VR face-to-face interaction study comparing stylised and photorealistic avatars. Stylised avatars produced significantly lower eeriness and significantly higher trustworthiness ratings, with no interaction between stylisation level and facial expression intensity.
}

\entry{Cihodaru-\c{S}tefanache \& Podina}{2025}{
A systematic review of 29 empirical studies on embodied conversational agents. Identifies consistent stylisation across all visual features as the most reliably recommended design strategy for avoiding uncanny valley responses, alongside accurate lip synchronisation and visual/voice realism consistency.
}

\entry{Weidner et al.}{2023}{
A systematic review of 72 studies on avatar and agent visualisation in AR and VR. Full-body, realistically rendered avatars generally produced stronger presence, embodiment, and task performance, but this advantage was task-dependent and sometimes reversed in game-like or educational scenarios.
}

\section*{Measurement instruments}

\entry{Bartneck et al.}{2009}{
Introduces the Godspeed Questionnaire Series, validated semantic differential scales across five perceptual dimensions: anthropomorphism, animacy, likeability, perceived intelligence, and perceived safety. Originally developed for human-robot interaction and since widely adapted for virtual avatar studies.
}

\entry{MacDorman \& Ishiguro}{2006}{
Advances the uncanny valley from a qualitative hypothesis into an empirically tractable research programme using androids as controllable experimental stimuli. Proposes semantic differential instruments to quantify eeriness and strangeness as measurable affective dimensions.
}

\entry{Ho \& MacDorman}{2010}{
Refines eeriness measurement into a two-subscale instrument that distinguishes perceived eeriness from perceived human-likeness. Finds the two dimensions are empirically orthogonal (correlation $r \approx .02$), meaning eeriness cannot be explained as simply the inverse of perceived humanness.
}

\end{document}
